% What's the problem?
% Why is it a problem? Research gap left by other approaches?
% Why is it important? Why care?
% What's the approach? How to solve the problem?
% What's the findings? How was it evaluated, what are the results, limitations, 
% what remains to be done?

% XXX Summary
\emph{Summary:}
\dots

% Motivation and intended learning outcomes
\emph{Avsedda lärandemål:}
Efter att ha genomfört detta kapitel ska du kunna:
\begin{itemize}
  \item Förklara vad en fil är och skillnaden mellan text- och
    binärfiler. % Lo10
  \item Öppna och stänga filer på rätt sätt med \mintinline{python}{open}
    och \mintinline{python}{with}, och förklara varför filen måste
    stängas. % Lo10
  \item Läsa data från en fil rad för rad, tolka raderna och lagra dem i
    lämpliga behållare. % Lo10, Lo09
  \item Skriva data från programmet till en fil i ett format som
    programmet senare kan läsa tillbaka. % Lo10
  \item Hantera fel vid filhantering, som att filen inte finns, så att
    användaren får en begriplig fråga i stället för en krasch. % Lo08
  \item Läsa och skriva vanliga filformat som CSV och JSON med hjälp av
    Pythons standardbibliotek och dess dokumentation. % Lo10, Lo06
\end{itemize}

% XXX Prerequisites
\emph{Prerequisites:}
\dots

% XXX Reading material
\emph{Reading:}
\dots
